\FrontHeading{Abstract}

People in Nepal who want skincare advice face three barriers. Dermatologists are concentrated in a few cities. Most AI skincare tools are trained on lighter-skinned populations. Many of the products those tools suggest cannot be bought locally. This project set out to build and evaluate \emph{SkinCare AI}, a mobile prototype that looks at a face photograph, reports visible cosmetic skin concerns and recommends products sold by Nepalese online retailers.

The artefact changed substantially during the project. The interim plan proposed a disease-and-skin-type classifier. An audit of more than 26,000 candidate images found widespread duplication, stock imagery, pre-generated augmentations and unusable skin-type labels. As a result, the scope was narrowed to five non-medical concerns: blemishes, dark spots, redness, visible pores and fine lines. Skin type is now reported by the user rather than guessed from a photo.

The final system has five parts:
\begin{itemize}
  \item an EfficientNetB0 multi-label model trained with masked loss on 1,200+ human-reviewed and FFHQ-Wrinkle images;
  \item a YuNet face-quality gate;
  \item a catalogue of 1,500+ products scraped from Jeevee, ForEveryNG and Oriflame;
  \item a transparent weighted recommender;
  \item a FastAPI backend with a React Native (Expo) client and an optional Qwen2.5 chat assistant.
\end{itemize}

Four model versions were compared on the same duplicate-group-aware held-out set of 470+ images. The deployed model, retrained after a second round of human label review, reached a pooled label accuracy of 76.3\%, a macro F1 of 0.789, a macro balanced accuracy of 0.670 and a macro ROC-AUC of 0.734, improving on the previous version on all four. Fine lines were detected reliably (ROC-AUC 0.920, balanced accuracy 0.833). Dark spots (ROC-AUC 0.580) and visible pores (0.597) remained close to chance. Class reweighting and extra melasma images were tested and rejected because they did not help.

The finished system was also tested directly. Eighteen of the 19 functional API tests passed; the chat assistant timed out under load. The face gate rejected all 300 out-of-distribution images. However, the gate also rejected 22\% of held-out skin close-ups, which showed that the gate and the model do not accept the same kinds of image. Usability was assessed through a heuristic evaluation and a participant survey run through Google Forms, which combined task questions with the System Usability Scale.

The project contributes an honest, leakage-aware baseline and a working, Nepal-specific pipeline from photo to purchasable product. Its main lesson is that data quality, not model choice, limits cosmetic skin analysis at this scale.
